\section{Recursive equilibrium}\label{sec:recursive}

\subsection{Uniqueness gives the minimal state}

\begin{theorem}\label{thm:minimal}
If sequential equilibrium is unique from every $(z,\mu)$, then there is a recursive equilibrium on
the minimal state $(z,\mu)$: the distribution next period is a function of today's shock, today's
distribution and tomorrow's shock, weakly continuous in the distribution, and saving is a continuous
function of today's state.
\end{theorem}
\formal{KSModel.isRecursive, continuous\_lawOfMotion}

\noindent
The continuation of a sequential equilibrium below a node is a sequential equilibrium from that
node's state, and the graph of the equilibrium correspondence is closed; with uniqueness the
equilibrium path is a continuous function of $(z,\mu)$ by the closed-graph theorem on a compact
space, and the law of motion is read off it. This is \citet[Proposition 2]{cao2020}.

\subsection{Enlarged states, in every economy}

\begin{theorem}\label{thm:valuestate}
In every economy there is a recursive equilibrium whose state is $(z,\mu,v)$, with $v$ the profile
of household values, and every such equilibrium generates a sequential equilibrium with the same
payoffs.
\end{theorem}
\formal{recursiveEquilibrium\_valueState, generates\_sequentialEquilibrium, payoffEquivalent}

\begin{theorem}\label{thm:history}
In every economy there is a recursive equilibrium whose aggregate state is the infinite history of
aggregate shocks: an assignment $\mu^*$ of a distribution to each history, consistent with optimal
saving at every history, with the history evolving exogenously. This holds in particular with
logarithmic utility and Cobb--Douglas prices.
\end{theorem}
\formal{exists\_seqStateEquilibrium, exists\_historyRecursiveEquilibrium,
logEconomy\_exists\_seqStateEquilibrium}

\noindent
The fixed point is taken in the space of assignments of distributions to infinite histories; no
initial distribution is needed. We call it the \emph{sequence-state equilibrium}. Under uniqueness
of sequential equilibrium it is Markov in $(z,\mu)$.

\begin{theorem}\label{thm:dgmm}
In every economy there is a stationary recursive equilibrium with an ergodic invariant distribution,
on the expectations-augmented state of \citet{dgmm1994}; it can be chosen conditionally spotless in
the sense of \citet[Definition 6]{cao2020}.
\end{theorem}
\formal{exists\_stationaryRecursiveEquilibrium, exists\_conditionallySpotlessEquilibrium}

\noindent
The proof formalises DGMM's Theorem 1.1 and Corollary 1.1 in the abstract (Kakutani--Fan--Glicksberg
for a correspondence of measures, disintegration, Krein--Milman for the ergodic choice) and applies it
to the compact set of equilibrium states.

\subsection{The minimal state without uniqueness}

\begin{theorem}\label{thm:monotone}
If saving rises with the path of aggregate capital, the greatest sequential equilibrium is a
consistent selection, and a recursive equilibrium on $(z,\mu)$ exists without uniqueness. Saving
rises with capital in complementarity economies (a constant wage, a return nondecreasing in capital,
relative risk aversion at most one), and an explicit such economy is given.
\end{theorem}
\formal{SavingMonotone.isRecursive, savingMonotone\_of, example\_exists\_consistentSelection}

\noindent
The condition fails for Cobb--Douglas prices, where higher capital lowers the cash on hand of
wealthy households.
